\documentclass[12pt]{article}
\usepackage[utf8]{inputenc}
\usepackage[T1]{fontenc}
\usepackage{hyperref}
\title{Papora-Hoanya Language Resources for K-12 Curriculum Development}
\author{Compiled by Claude Sonnet 5.5 (claude-sonnet-5-5) for Andr\'as Kornai}
\date{}
\begin{document}
\maketitle
\section{Introduction}
Papora-Hoanya (ISO 639-3 ppu) is an extinct Formosan (Austronesian) language of central-western Taiwan. The Papora are among the Pingpu groups seeking official recognition. No dedicated Papora-Hoanya revival project, dictionary or teaching programme was found in the searches.
\section{Primary Educational Resources}
None found.
\section{Linguistic and Reference Materials}
\begin{itemize}
\item Wikipedia, Papora-Hoanya language: \url{https://en.wikipedia.org/wiki/Papora-Hoanya_language}
\item Journal article (UEA repository) surfaced by the search; its relevance to Papora-Hoanya was not verified: \url{https://ueaeprints.uea.ac.uk/99055/1/ijts-article-10.1163-24688800-20251427.pdf}
\end{itemize}
\section{Religious and Community Resources}
\begin{itemize}
\item The China Project, on unrecognized Taiwanese indigenous tribes reviving languages (Pingpu context): \url{https://thechinaproject.com/2023/06/05/taiwans-unrecognized-indigenous-tribes-are-reviving-dead-languages-to-achieve-recognition/}
\end{itemize}
\section{Digital and Audio Resources}
None found.
\section{Curriculum Development Context}
Comparable Pingpu revival work, such as Siraya revitalisation from historical texts, may offer a model: \url{https://taiwaninsight.org/2026/03/27/language-revitalisation-using-historical-texts-the-case-of-siraya/}. Background on national policy: \url{https://scholarspace.manoa.hawaii.edu/items/847ac3de-a2cd-44a0-afdf-d83506ffc3b8/full}. These cover recognized groups; no programme targeting Papora-Hoanya was found.
\end{document}